\documentclass[11pt]{article}
\usepackage[margin=1.05in]{geometry}
\usepackage{amsmath,amssymb,amsthm,booktabs,microtype,array,xcolor}
\usepackage[colorlinks=true,linkcolor=blue!55!black,urlcolor=blue!55!black]{hyperref}
\newtheorem{definition}{Definition}[section]
\newtheorem{theorem}[definition]{Theorem}
\newtheorem{conjecture}[definition]{Conjecture}
\newtheorem{example}[definition]{Example}
\newtheorem{counterexample}[definition]{Counterexample}
\newtheorem{finding}[definition]{Finding}
\newcommand{\lean}[1]{\texttt{#1}}
\newcommand{\Ev}{\mathsf{Ev}}
\newcommand{\seen}{\mathsf{seen}}
\newcommand{\merge}{\mathsf{merge}}
\newcommand{\live}{\mathsf{live}}
\newcommand{\tok}{\mathsf{tok}}
\newcommand{\pos}{\mathsf{pos}}
\newcommand{\cells}{\mathsf{cells}}
\newcommand{\ranges}{\mathsf{ranges}}
\newcommand{\mvr}{\mathsf{mvr}}
\newcommand{\Row}{\mathsf{row}}
\newcommand{\Col}{\mathsf{col}}
\sloppy
\title{A Materialised Three-Way Merge for AegisSheet\\
  \large Exploratory research note}
\author{FP Launchpad}
\date{2 September 2026}
\begin{document}
\maketitle

\begin{abstract}
The verified AegisSheet model stores the whole event history and merges by
union; each event carries the timestamp set of its causal past, so state is
quadratic in history length. This note asks whether the ancestor argument of
a three-way merge can replace that metadata. We define a materialised state
with a componentwise three-way merge and test it differentially against the
existing model on random version DAGs. No counterexample appears in 3000
executions when a removal names the tokens it kills; clearing tokens instead
converges but breaks commutation. Three hand-derived witnesses, checked by
the executable reference, show that the last position of a removed row or
column cannot be dropped while a concurrent revival is possible, so the
retention residue is one register per axis identifier until removal is
causally stable. The measured state is linear rather than quadratic in
history. None of the new claims is machine-checked; the mechanization
target is a conditioned Join for the token and version components composed
with the register, plus an observational-equivalence theorem against the
union model.
\end{abstract}

\section{Goal, hypotheses, and status}

\paragraph{Goal.} Re-encode AegisSheet as a genuine three-way MRDT with a
materialised state so that the lowest-common-ancestor argument of $\merge$
replaces the per-event causal set, the append-only event log, and, where
possible, the roster-acknowledged purge protocol, while preserving every
observation of the current model.

\paragraph{Hypotheses and falsifiers.}
\begin{description}
\item[H1] Cells and axes forget. A state of live keep tokens, latest-position
registers, active cell versions, and active range versions, merged
componentwise by the rule $\mvr(l,a,b)=(l\cap a\cap b)\cup(a\setminus l)
\cup(b\setminus l)$ for sets and last-writer-wins for registers, yields the
same observation as the union model at every reachable version.
Falsifier: a version DAG on which the observations differ, two merge
topologies over one event set that disagree, or a merge that differs from
replaying the union.
\item[H2] The latest position of a removed axis identifier can be dropped
once no live range names it. Falsifier: two histories with identical live
sheets whose continuations require different positions for the identifier.
\item[H3] Purge dissolves into ordinary version deletion once state is
materialised. Falsifier: a late concurrent write that revives content the
paper requires to stay purged.
\end{description}

\paragraph{Status.} Table~\ref{tab:status} uses the evidence vocabulary
conjectured, refuted, machine-checked, validated, measured, specified, and
staged. None of the candidate claims in this note is machine-checked; the
reference definitions and the baseline theorems cited from the repository
are. The Python reference is validated against the Lean model only on the
fixtures listed in Section~\ref{sec:controls}.

\begin{table}[h]
\centering\small
\begin{tabular}{@{}p{0.30\textwidth}p{0.16\textwidth}p{0.46\textwidth}@{}}
\toprule
Claim & Status & Evidence \\
\midrule
H1, named removal & validated on honest executions; conjectured in general & no counterexample in 3000 executions under the D1 issuance rule, 118{,}474 versions, 21{,}571 registered-ancestor and 834 virtual-base merges, three seeds; all 96{,}069 generated events pass a transcribed \lean{applicableB} after erasing $\mathit{kills}$; Lean proof owed \\
H1, clearing removal & validated for observation; refuted for replay & converges and matches the reference; 121 of 3000 executions have a merge that differs from timestamp-order replay \\
H2 & refuted & three hand-derived witnesses checked by the reference (Section~\ref{sec:h2}); 11 to 55 failing executions per 1000 in the harness \\
H3 & staged & purge is not modelled by the harness \\
State growth & measured & Table~\ref{tab:growth}: reference quadratic, materialised linear \\
Undo revives a later removal & finding, reproduced; excluded by decision D1 & Section~\ref{sec:undo}; legacy behaviour of the current model \\
\bottomrule
\end{tabular}
\caption{Status of each claim.}
\label{tab:status}
\end{table}

\paragraph{Trusted boundary.} The reference semantics is the executable
Lean model in \lean{Sal/MRDTs/Instances/AegisSheet.lean}, which the
repository validates against Tables~3 and~4 of the AegisSheet paper. The
Python transcription of that model is trusted only to the extent of the
fixture check. The candidate design is specified here and implemented in
Python; it has no Lean declaration.

\section{The reference model}
\label{sec:reference}

The displays below transcribe the Lean definitions named in each definition's title.
Purge markers are omitted throughout (labelled compression); the harness does
not generate them.

\begin{definition}[events and state; \lean{SheetOp}, \lean{D}]
An event is $e=(t,r,\langle\seen,\mathit{cmd}\rangle)$ with Lamport
timestamp $t$, replica $r$, a finite set $\seen$ of timestamps, and a command
that is either a direct action or $\mathsf{undo}(t',\mathit{inv})$ carrying
an inverse action. The effect $\mathsf{action}(e)$ is the direct action or
$\mathit{inv}$. Actions are
\[
\begin{aligned}
&\mathsf{axis}(\mathit{kind},\alpha,\mathit{id},\mathit{before},\mathit{after})
 &&\alpha\in\{\Row,\Col\},\ \mathit{kind}\in\{\mathsf{insert},\mathsf{move},\mathsf{remove},\mathsf{restore}\},\\
&\mathsf{cell}(\rho,\gamma,\mathit{before},\mathit{after},\mathit{overwrites})
 &&\mathit{before},\mathit{after}\text{ finite value sets},\ \mathit{overwrites}\text{ timestamps},\\
&\mathsf{range}(\mathit{id},\mathit{before},\mathit{after},\mathit{overwrites})
 &&\mathit{before},\mathit{after}\in\mathsf{Option}(\mathit{RangeSpec}).
\end{aligned}
\]
The MRDT state is a finite set $E$ of events; $\mathsf{do}(E,e)=E\cup\{e\}$;
$\merge(\_,a,b)=a\cup b$. The ancestor argument is unused.
\end{definition}

\begin{definition}[axis liveness; \lean{keepsAxis}, \lean{removesAxis}, \lean{axisLive}]
An event $e$ \emph{keeps} $(\alpha,\mathit{id})$ when it is an axis update of
that identifier with $\mathit{after}\neq\bot$, or a cell update whose
$\alpha$-coordinate is $\mathit{id}$. It \emph{removes} $(\alpha,\mathit{id})$
when it is an axis update of that identifier with $\mathit{after}=\bot$.
\[
\begin{aligned}
\mathsf{keepTimes}_E(\alpha,\mathit{id})&=\{t(e)\mid e\in E,\ e\text{ keeps }(\alpha,\mathit{id})\},\\
\mathsf{removed}_E(\alpha,\mathit{id},\tau)&\Longleftrightarrow
 \exists e\in E.\ e\text{ removes }(\alpha,\mathit{id})\land\tau\in\seen(e),\\
\live_E(\alpha,\mathit{id})&\Longleftrightarrow
 \mathsf{known}_E(\alpha,\mathit{id})\land
 \exists\tau\in\mathsf{keepTimes}_E(\alpha,\mathit{id}).\ \neg\mathsf{removed}_E(\alpha,\mathit{id},\tau),
\end{aligned}
\]
where $\mathsf{known}_E$ holds when some axis update of the identifier is in
$E$. A removal defeats exactly the keep tokens it observed; a keep it did not
observe survives it.
\end{definition}

\begin{definition}[positions, cells, ranges; \lean{axisPositions}, \lean{cellValues}, \lean{rangeValues}]
A candidate for $(\alpha,\mathit{id})$ is an axis update with
$\mathit{after}\neq\bot$; $\mathsf{positions}_E(\alpha,\mathit{id})$ is the
set of $\mathit{after}$ values of candidates with no later candidate, a
singleton or empty. A cell version $e$ at $(\rho,\gamma)$ is overwritten when
some cell update in $E$ lists $t(e)$ in its $\mathit{overwrites}$;
$\cells_E(\rho,\gamma)$ is the union of $\mathit{after}$ over non-overwritten
versions when both axes are live, and $\emptyset$ otherwise. Range versions
are analogous, and $\ranges_E(\mathit{id})$ collects their non-$\bot$
$\mathit{after}$ specifications.
\end{definition}

\begin{definition}[range resolution; \lean{resolveRange}]
For a specification $(\rho_1,\rho_2,\gamma_1,\gamma_2)$, a live endpoint
resolves to itself. A dead first endpoint resolves to the live identifier at
the least live position greater than the endpoint's last position; a dead
last endpoint to the live identifier at the greatest live position less than
it. Resolution fails when an endpoint has no position or the resolved
rectangle is inverted.
\end{definition}

\begin{definition}[issuance; \lean{applicableB}, \lean{validUndo}, \lean{inverseFor}]
An event is issuable at $E$ when $t(e)\notin\mathsf{times}(E)$,
$\seen(e)=\mathsf{times}(E)$, $t(e)$ exceeds every timestamp in $E$, and its
before-image and $\mathit{overwrites}$ list equal the current values at $E$
(compressed: the per-kind clauses of \lean{directApplicable}). An undo is
issuable when its target is in $\seen(e)$, was issued by the same replica,
and the carried inverse equals $\mathsf{inverseFor}$ of the target, which
swaps before and after and sets $\mathit{overwrites}=\{t(\text{target})\}$
for cells and ranges. No clause requires the axes of an undone cell write
to be live.
\end{definition}

The clause $\seen(e)=\mathsf{times}(E)$ is the cost this note attacks: every
event stores the timestamp set of its entire causal past.

\section{The candidate design}
\label{sec:design}

\begin{definition}[materialised state]
\[
\begin{array}{ll}
\mathsf{known}:\{\Row,\Col\}\to\mathcal P(\mathit{Id}) & \text{grow-only identifier sets},\\
\tok:(\{\Row,\Col\}\times\mathit{Id})\to\mathcal P(\mathit{Ts}) & \text{live keep tokens},\\
\pos:(\{\Row,\Col\}\times\mathit{Id})\rightharpoonup\mathit{Ts}\times\mathit{Position} & \text{latest position register},\\
\cells:(\mathit{Id}\times\mathit{Id})\to\mathcal P(\mathit{Ts}\times\mathcal P(\mathit{Value})) & \text{active cell versions},\\
\ranges:\mathit{Id}\to\mathcal P(\mathit{Ts}\times\mathsf{Option}(\mathit{RangeSpec})) & \text{active range versions}.
\end{array}
\]
A removal carries $\mathit{kills}$, the live tokens of its identifier at
issuance (the \emph{named} variant); the \emph{clearing} variant carries
nothing and empties the token set.
\end{definition}

\begin{definition}[update]
\[
\begin{aligned}
\mathsf{axis}(\_,\alpha,\mathit{id},\_,p\neq\bot)&:\ \tok(\alpha,\mathit{id})\mathrel{+}=\{t\},\quad
 \pos(\alpha,\mathit{id}):=\max_{\mathrm{ts}}(\pos(\alpha,\mathit{id}),(t,p)),\quad \mathsf{known}\mathrel{+}=\mathit{id};\\
\mathsf{axis}(\_,\alpha,\mathit{id},\_,\bot)&:\ \tok(\alpha,\mathit{id})\mathrel{-}=\mathit{kills}
 \quad(\text{clearing: }\tok(\alpha,\mathit{id}):=\emptyset);\\
\mathsf{cell}(\rho,\gamma,\_,v,\mathit{ov})&:\ \tok(\Row,\rho)\mathrel{+}=\{t\},\ \tok(\Col,\gamma)\mathrel{+}=\{t\},\\
 &\quad\cells(\rho,\gamma):=\{x\in\cells(\rho,\gamma)\mid \mathrm{ts}(x)\notin\mathit{ov}\}\cup\{(t,v)\};\\
\mathsf{range}(\mathit{id},\_,s,\mathit{ov})&:\ \ranges(\mathit{id}):=\{x\in\ranges(\mathit{id})\mid\mathrm{ts}(x)\notin\mathit{ov}\}\\
 &\qquad\cup\{(t,s)\}.
\end{aligned}
\]
\end{definition}

\begin{definition}[merge]
For $\mvr(l,a,b)=(l\cap a\cap b)\cup(a\setminus l)\cup(b\setminus l)$,
\[
\begin{aligned}
\merge(l,a,b).\mathsf{known}&=l.\mathsf{known}\cup a.\mathsf{known}\cup b.\mathsf{known},\\
\merge(l,a,b).\tok(k)&=\mvr(l.\tok(k),a.\tok(k),b.\tok(k)),\\
\merge(l,a,b).\pos(k)&=\max_{\mathrm{ts}}\{x.\pos(k)\mid x\in\{l,a,b\},\ k\in\mathrm{dom}(x.\pos)\},\\
\merge(l,a,b).\cells(k)&=\mvr(\ldots),\qquad
\merge(l,a,b).\ranges(k)=\mvr(\ldots).
\end{aligned}
\]
The observation is computed as in Section~\ref{sec:reference} with
$\live(\alpha,\mathit{id})\Leftrightarrow\mathit{id}\in\mathsf{known}(\alpha)
\land\tok(\alpha,\mathit{id})\neq\emptyset$ and the register in place of the
latest candidate.
\end{definition}

The design has no cross-component rule. The reference's update-wins policy
(a concurrent write or move defeats a removal) is already token-based: the
write's token is not in the ancestor, so $\mvr$ keeps it. Every component is
an observed-remove set or a last-writer-wins register.

\section{H1: differential campaign}
\label{sec:h1}

\begin{conjecture}[H1]
For every legal execution of the version-DAG semantics with the
canonical virtual merge base, the materialised design with named removal has
the same observation as the union model at every version.
\end{conjecture}

\paragraph{Method.} The harness (\lean{model.py}) runs the Python
reference and each materialised variant in lockstep on random version DAGs.
Operations are generated from the reference state so that before-images,
overwrite lists, and killed-token lists are those the issuance predicate
demands; identifiers and timestamps are fresh by construction. Two to four
replicas perform one to three operations each per round for three to eight
rounds; a merge between two heads occurs with probability $0.7$ per round;
forced convergence rounds then merge all heads along different topologies.
When a recorded version's event set equals the heads' intersection it is the
ancestor; otherwise the ancestor is the fold of the intersection in
timestamp order, which is the canonical virtual merge base. At every version
the harness compares observations, compares observations of versions
reached by different topologies over the same event set, and compares the
merge result with the timestamp-order replay of the union. Every generated
event, after erasing $\mathit{kills}$, is checked against a transcription of
\lean{applicableB} and \lean{clockedB}; all 96{,}069 events of the three
seeds pass, so the executions are honest in the Lean sense and not only in
the generator's. The generator applies decision D1
(Section~\ref{sec:undo}): an inverse cell write is issued only while its
axes are live.

\paragraph{Result.} Table~\ref{tab:campaign} reports three seeds of 1000
executions. The named-removal design has no failing execution. The clearing
variant matches the reference and converges but fails the replay comparison:
a removal that clears all tokens and a concurrent keep do not commute, so
replay order matters and a Lean proof would need the arbitration arm of the
replay policy. The named variant makes every \emph{concurrent} pair commute.
Causally ordered pairs still do not: a keep followed by the removal that
names its token leaves the token absent, the reverse order leaves it
present, and the same holds for an overwrite and the version it names. So
universal commutation fails and the all-commuting proof route is
unavailable; Section~\ref{sec:mech} gives the conditioned route instead.

\begin{table}[h]
\centering\small
\begin{tabular}{@{}lrrrl@{}}
\toprule
Design & seed 1 & seed 2 & seed 3 & first failure class \\
\midrule
keep registers, named removal & 0 & 0 & 0 & \\
keep registers, clearing removal & 35 & 31 & 55 & merge $\neq$ replay \\
drop unreferenced dead registers, named & 16 & 11 & 11 & position of revived axis \\
drop all dead registers, named & 9 & 6 & 9 & range resolution \\
ancestor ignored (control) & 179 & 156 & 141 & cells, ranges, rows \\
eager range re-anchoring (control) & 646 & 675 & 684 & ranges \\
\bottomrule
\end{tabular}
\caption{Failing executions out of 1000 per seed under the D1 issuance
rule. Each seed produced about 39{,}500 versions, 32{,}000 operations,
7{,}200 merges at a registered ancestor, and 280 merges at a virtual base.
Under the legacy undo rule the counts are within a few executions of these
and the named design is likewise clean.}
\label{tab:campaign}
\end{table}

\section{H2: the last position of a dead axis is load-bearing}
\label{sec:h2}

Each witness is hand-derived from the definitions of
Section~\ref{sec:reference} and checked by evaluating the Python reference;
the checks are the functions \lean{h2\_witnesses} and \lean{h2\_pair3}.
Identifiers: rows $r_0,r_1,r_2$ at positions $10,20,30$, column $c_0$.

\begin{counterexample}[range resolution]
Let $L$ insert $r_0,r_1,r_2,c_0$ and a range $(r_0,r_2,c_0,c_0)$. History
$W_1$ removes $r_0$. History $W_2$ moves $r_0$ to position $25$ and then
removes it. Both leave live rows $\{r_1,r_2\}$ at unchanged positions and
the same range specification. The reference resolves the range to
$(r_1,r_2,c_0,c_0)$ after $W_1$ and to $(r_2,r_2,c_0,c_0)$ after $W_2$. A
state that forgets $r_0$'s last position cannot give both answers.
\end{counterexample}

\begin{counterexample}[eager re-anchoring]
From the same $L$, branch $A$ removes $r_0$ and branch $B$ concurrently
inserts $r_3$ at position $15$. The reference resolves the merged range to
$(r_3,r_2,c_0,c_0)$. Rewriting the range to $r_0$'s successor at removal
time gives $(r_1,r_2,c_0,c_0)$. Lazy resolution over the dead position is
part of the specification, not an implementation choice.
\end{counterexample}

\begin{counterexample}[update-wins revival]
Let $L$ insert $r_0$ at $10$ and $c_0$, and write cell $(r_0,c_0)$. Branch
$A_1$ moves $r_0$ to $52$ and then removes it; branch $A_2$ only removes it.
Both branch states have identical live sheets: $r_0$ is dead. Branch $B$
concurrently writes $(r_0,c_0)$. In both merges the write's token revives
$r_0$; the reference places it at $52$ after $A_1$ and at $10$ after $A_2$.
The merge receives identical $L$ and $S_B$, so if $S_{A_1}=S_{A_2}$ it
cannot answer both. The information lives only on the removing branch.
\end{counterexample}

The third counterexample does not involve ranges or undo, and it survives
the D1 issuance rule (11 to 16 failing executions of 1000 per seed for the
register-dropping variant). The
retention residue is therefore: one position register per axis identifier
that has ever been known, retained until the removal is causally stable,
that is, until no replica can still issue an operation concurrent with it.
Under update-wins this is unavoidable; under a remove-wins policy the
register could be dropped at removal. Causal stability suffices only if undo
cannot revive a removed axis: under the current model
(Section~\ref{sec:undo}) an undo issued \emph{after} observing the removal
still revives it, and the register must then be kept as long as any replica
can undo a write into the row. Decision D1 adopts the no-revival rule, so
causal stability is the retirement condition the purge experiment must test.

\section{Finding: undo revives a causally later removal}
\label{sec:undo}

\begin{finding}
In the current model, with row $r_0$ and column $c_0$ and a write at $t=3$,
removing $r_0$ at $t=4$ and then undoing the $t=3$ write at $t=5$ is a legal
history on one replica: the undo's target is observed and was issued by the
same replica, and no issuance clause requires the axes to be live. The undo
is a cell event and therefore a keep token for $r_0$ with a timestamp the
removal did not observe. The row becomes live again at position $10$ with an
empty cell. Reproduced by \lean{undo\_revival\_witness}.
\end{finding}

Table~4 of the paper covers undo against a \emph{concurrent} removal; it
does not say what undoing an edit should do to a row removed causally later.
The behaviour is a consequence of representing undo as a cell event. If it is
unintended, the fix is an issuance clause (undo of a cell write requires live
axes) or a non-keeping inverse; either changes the public sequential
specification.

\paragraph{Decision D1.} No revival, as an issuance rule: an inverse cell
write is legal only while its row and column are live. The legacy revival
behaviour stays documented by this finding and its SPOT; the equivalence
theorem of Section~\ref{sec:mech} is stated only for honest executions under
the new rule.

\section{H3: purge}

The harness does not generate purge markers, so H3 is staged. Two
observations shape the next experiment. First, in the materialised design a
purge is an ordinary deletion of cell versions at dead coordinates, and the
$\mvr$ rule already keeps a concurrent write's version because it is absent
from the ancestor, which is the paper's stated intent that a late change
restores the axis but not the purged contents. Second, the roster
acknowledgement of the current protocol is the causal-stability condition
that Section~\ref{sec:h2} identifies as the moment the dead register can be
dropped. The conjecture for the next step is that the acknowledgement
protocol is needed for the $O(1)$ register, not for the cell payloads.

\section{Measured: state size}

\begin{table}[h]
\centering\small
\begin{tabular}{@{}rrrr@{}}
\toprule
Rounds & Operations & Reference (timestamps stored) & Materialised, named, keep \\
\midrule
4 & 23 & 131 & 21 \\
8 & 46 & 593 & 44 \\
16 & 92 & 2{,}902 & 84 \\
32 & 185 & 13{,}835 & 180 \\
\bottomrule
\end{tabular}
\caption{Mean over 20 executions with 3 replicas, seed 1, at the final
version of replica 0. The reference count is one per event plus its
$\seen$ set and overwrite list; the materialised count is tokens, registers,
and versions.}
\label{tab:growth}
\end{table}

Across the measured range, doubling the history roughly quadruples the
reference count and doubles the materialised count. This is a
timestamp-count measurement consistent with quadratic versus linear scaling,
not an asymptotic result. The gap at 185 operations is a factor of 77. The
counts are taken in the Python models, not bytes of any runtime.

\section{Controls}
\label{sec:controls}

\paragraph{Positive control.} The Python reference reproduces twenty
hand-derived expectations copied from the Lean SPOTs, and five issuance
fixtures (honest and dishonest edits, own and remote undo, undo of insert): update defeats
concurrent removal, later removal wins, concurrent writes keep both values,
selective undo keeps the remote value, move defeats removal, later move wins,
insert survives adjacent removal, undo of remove restores the position, undo
of insert removes only the inserted identifier, the two Table~4 move-undo
cases, and seven Figure~1 range cases (\lean{check\_fixtures}).

\paragraph{Negative controls.} Two deliberately wrong designs fail on every
seed (Table~\ref{tab:campaign}): ignoring the ancestor, which reintroduces
removed versions, and eager range re-anchoring. A harness that could not
distinguish these would carry no evidence.

\section{Mechanization correspondence}
\label{sec:mech}

\begin{center}\small
\begin{tabular}{@{}p{0.34\textwidth}p{0.30\textwidth}p{0.28\textwidth}@{}}
\toprule
Display & Lean declaration & Status \\
\midrule
Reference events, state, liveness, positions, cells, ranges, resolution, issuance & \lean{Instances.AegisSheet}: \lean{D}, \lean{axisLive}, \lean{axisPositions}, \lean{cellValues}, \lean{rangeValues}, \lean{resolveRange}, \lean{applicableB} & machine-checked definitions; transcription validated on 20 fixtures \\
Materialised state, update, merge & none & specified; Python implementation \\
Conjecture H1 & none & validated on the campaign scope \\
Counterexamples 5.1 to 5.3, Finding 6.1 & none & hand-derived, checked by the Python reference; to be pinned as Lean SPOTs \\
\bottomrule
\end{tabular}
\end{center}

\paragraph{Lean plan.} The materialised design is a product of
observed-remove sets and one last-writer-wins register. Because causally
ordered pairs do not commute, the all-commuting route is closed. The token
and version components need \lean{JoinAt} under an honesty predicate on the
replay context (every removal's $\mathit{kills}$ and every overwrite's
$\mathit{overwrites}$ lie in the issuer's causal past), discharged from
issuance with \lean{IssuanceEstablishes} as \lean{Queue} and \lean{EmbedRGA}
do, and composed with \lean{joinAt\_prod}; the register alone is
all-commuting. A \lean{VerifiedMRDT} for the new package proves correctness
against its own sequential specification only. H1 itself needs a separate
theorem: an erasure dropping $\mathit{kills}$ that maps honest executions of
the new package to honest executions of the union model, a representation
relation between materialised states and event sets preserved by
$\mathsf{do}$ and $\merge$, and equality of observations under it. The
sequential certificate needs a new representation relation;
\lean{no\_view\_only\_step} shows the visible sheet is insufficient, and
\lean{row\_tokens\_distinguish\_future},
\lean{cell\_versions\_distinguish\_future}, and
\lean{range\_versions\_distinguish\_future} fix that tokens, cell version
identities, and range version identities must be retained, which the
materialised state does.

\section{Open obligations}

\begin{enumerate}
\item Decision D1 (no revival) is taken; carry it into the Lean issuance
predicate of the new package and keep the legacy SPOT.
\item Model purge in the harness under the chosen policy and test H3: purge
as version deletion, concurrent writes surviving, and the register-dropping
rule, with the undo-after-stability case as a control.
\item Pin Counterexamples 5.1 to 5.3 and Finding 6.1 as Lean SPOTs against
the existing model, with PASS and FAIL companions.
\item Port the named-removal design to Lean as a second package beside the
current one; prove \lean{JoinOn} under the honesty predicate, the sequential
certificate, and the cross-model observational-equivalence theorem; keep the
union model until the runtime measurement is in.
\item Extend the harness to merges between non-head versions and to undo of
undo; report the generator's coverage.
\end{enumerate}

\end{document}
